\documentclass[11pt]{article}

\usepackage[letterpaper,margin=0.9in]{geometry}
\usepackage{amsmath,amssymb,amsfonts}
\usepackage{array,booktabs,longtable,tabularx}
\usepackage{enumitem}
\usepackage{titlesec}
\usepackage{fancyhdr}
\usepackage{lastpage}
\usepackage{hyperref}
\usepackage{xcolor}
\usepackage{graphicx}

\setlength{\headheight}{14pt}
\hypersetup{
  colorlinks=true,
  linkcolor=black,
  urlcolor=black,
  citecolor=black
}

\pagestyle{fancy}
\fancyhf{}
\lhead{Negative-Energy Modalities}
\rhead{Laboratory QFT Claim Grammar}
\cfoot{\thepage\ of \pageref{LastPage}}

\titleformat{\section}{\large\bfseries\uppercase}{\thesection.}{0.5em}{}
\titleformat{\subsection}{\normalsize\bfseries}{\thesubsection}{0.5em}{}
\setlist[itemize]{nosep,leftmargin=1.2em}
\setlist[enumerate]{nosep,leftmargin=1.4em}
\renewcommand{\arraystretch}{1.16}

\newcommand{\Tab}{Table}
\newcommand{\Tmunu}{T_{\mu\nu}}
\newcommand{\Tabmunu}{T_{ab}}
\newcommand{\Tnn}{T_{ab}k^ak^b}
\newcommand{\Ttt}{T_{00}}
\newcommand{\sbr}{\mathrm{S/B}}
\newcommand{\um}{\,\mu\mathrm{m}}

\title{\bfseries Negative-Energy Modalities in Laboratory Quantum Field Systems:\\
Boundary, State, Tensor, and Readout Claims}
\author{
Daniel S. Goldman\\
Independent researcher\\
\href{https://orcid.org/0000-0003-2835-3521}{ORCID: 0000-0003-2835-3521}
}
\date{\today}

\begin{document}
\maketitle

\begin{abstract}
Laboratory discussions of negative energy often join several distinct
physical objects under one phrase. A Casimir boundary condition, a prepared
nonclassical field state, a renormalized stress-tensor component, an averaged
energy-condition quantity, and an instrument response can each carry
negative-energy language while licensing different conclusions. This paper
sets out a compact claim grammar for those cases. The grammar separates four
modalities: boundary-conditioned response, state-conditioned response,
tensor/averaging relevance, and readout-conditioned observability.

The boundary-conditioned modality concerns vacuum response organized by
material or geometric constraints, as in Casimir systems and finite
worldline-numerics geometries. The state-conditioned modality concerns
negative stress-energy created or selected by field-state preparation,
measurement, squeezing, or quantum-energy-teleportation-type protocols. The
tensor/averaging modality concerns the stress-energy object itself: which
component is being discussed, how it is renormalized or smeared, which
averaged condition it is compared against, and what source role it can play.
The readout-conditioned modality concerns the apparatus channel: force,
pressure, force gradient, resonant frequency, phase, loss, current, voltage,
or detector response, together with the ordinary material and electromagnetic
controls that define the measurement.

This vocabulary turns negative-energy claims into transferable but
well-labeled statements. A boundary morphology can motivate a readout
schedule; a state-preparation protocol can supply a local stress-energy
deficit; a tensor calculation can compare the result with energy-condition or
source requirements; a precision instrument can test whether the corresponding
coefficient appears in measured data. The White et al. sphere--cylinder
Casimir result is used as a worked boundary-conditioned example: its
shell-like morphology, source-scale accounting, and physical readout routes
occupy different modalities that can be evaluated together while preserving
their meanings. The resulting framework is intended as a central reference for
a paper pair: one companion audit assigns the White geometry to the boundary
and tensor/source layers, and one companion readout paper asks how the
shell-coded object could enter pressure, cavity, and superconducting
measurement ensembles.
\end{abstract}

\section{Purpose}

Negative energy is a real and useful term in quantum field theory. It is also
a term with several laboratory meanings. The Casimir effect supplies boundary
conditioned vacuum response. Squeezed and measurement-conditioned states
supply state-dependent local stress-energy structures. Quantum energy
inequalities and averaged energy conditions supply duration, smearing, and
source constraints. Precision instruments supply force, pressure, frequency,
phase, loss, current, voltage, and detector observables. Each of those
statements belongs to a different physical modality.

The purpose of this paper is to make those modalities explicit. The goal is a
positive claim grammar: a way to say exactly what has been shown, what kind of
object carries the sign, which instrument can test it, and which additional
evidence promotes it into a stronger claim. This is especially useful for
finite Casimir geometries and exotic-spacetime-adjacent discussions, where a
boundary-generated negative distribution can be visually suggestive while
stress-tensor scale, averaged constraints, and apparatus readout remain
separate questions.

The motivating case is the White, Vera, Han, Bruccoleri, and MacArthur
sphere--cylinder geometry \cite{White2021}. In that work, worldline numerics
applied to custom Casimir cavities produced a shell-like negative energy
density distribution that was compared with the Alcubierre metric source
shape. That comparison intersects several mature parts of the field:
worldline methods for Casimir geometries \cite{Schafer2015}, energy
conditions and quantum energy inequalities \cite{Ford2009,KontouSanders2020},
precision Casimir metrology \cite{Bimonte2021,Stange2019}, and emerging
superconducting Casimir sensing \cite{Xu2026}. The value of the example is
that it asks all modality questions at once.

The central lesson is simple. A laboratory negative-energy claim becomes
clear when it identifies the modality first. Boundary-conditioned statements
describe how geometry organizes a vacuum response. State-conditioned
statements describe how preparation or measurement changes a field state.
Tensor/averaging statements describe the stress-energy object and its
constraints. Readout-conditioned statements describe what the apparatus can
recover against ordinary backgrounds. Those statements can support each other,
and their interfaces are where the strongest experimental programs live.

\section{Four Modalities}

A modality is the physical route by which a negative-energy statement obtains
meaning. It specifies the preparation or constraint, the stress-energy object,
the spatial and temporal support, the compensation structure, the readout
channel, and the inference licensed by the data.
\Tab~\ref{tab:modalities} gives the working classification.

\begin{table}[h]
\centering
\small
\begin{tabularx}{\textwidth}{p{0.18\textwidth}p{0.26\textwidth}X}
\toprule
Modality & Primary object & Licensed statement \\
\midrule
Boundary-conditioned &
Vacuum response fixed by material boundaries, geometry, and field boundary
conditions &
A finite apparatus geometry organizes Casimir energy, pressure, stress, or a
proxy density into a spatial pattern. \\
\addlinespace
State-conditioned &
Stress-energy in a prepared, post-selected, squeezed, measured, or
feed-forward-conditioned quantum field state &
A field state carries a local or smeared stress-energy deficit with a defined
preparation record and compensation structure. \\
\addlinespace
Tensor/averaging &
Renormalized stress tensor, null or timelike contraction, smeared integral, or
energy-condition quantity &
The signed object is tied to a specified component, averaging window, source
comparison, and quantum-inequality or averaged-condition context. \\
\addlinespace
Readout-conditioned &
Measured apparatus response: force, pressure, gradient, frequency, phase,
loss, current, voltage, transition, or detector signal &
A laboratory observable contains a recoverable coefficient with declared
controls, nuisance templates, calibration, and false-positive logic. \\
\bottomrule
\end{tabularx}
\caption{Four laboratory modalities for negative-energy claims. The
modalities can be combined, but each licenses a different kind of conclusion.}
\label{tab:modalities}
\end{table}

The same experiment can occupy several modalities. A structured Casimir cavity
can be boundary-conditioned in its source of vacuum response, tensor/averaging
conditioned in its comparison with \(\Tmunu\), and readout-conditioned in the
force or resonator signal used to measure it. A quantum-energy-teleportation
protocol can be state-conditioned in its preparation, tensor/averaging
conditioned in its local deficit, and readout-conditioned in the detector or
work-extraction channel used to score the event \cite{FunaiMartinMartinez2017}.

The classification is therefore an interface map rather than a taxonomy wall.
It helps a paper state which object is being advanced and which object is
being reserved for a companion calculation or experiment.

\section{Boundary-Conditioned Response}

The boundary-conditioned modality begins with a material or geometric
constraint on quantum fields. Parallel plates, sphere--plate systems,
microstructured surfaces, deep trenches, nanomembranes, and superconducting
microplates all fit this class. The signed object is the boundary-conditioned
change in vacuum energy, stress, pressure, force, or a calculational proxy for
those quantities. The physical control is geometry and material response.

This modality is the native home of Casimir physics. Precision measurements
typically recover force, force gradient, pressure, or resonator response, with
surface potentials, roughness, edge effects, finite conductivity, temperature,
alignment, and material optical data treated as part of the measurement model
\cite{Bimonte2021,Behunin2011,Garrett2015}. MEMS/NEMS and nanomembrane
platforms make the same point experimentally: the observable is an apparatus
response whose interpretation depends on careful material and electrostatic
accounting \cite{Decca2011,GarciaSanchez2012,Stange2019}.

Worldline numerics adds a valuable finite-geometry language to this modality.
It can evaluate Casimir energies and, in stress-tensor extensions, components
of the fluctuation-induced energy-momentum tensor in arbitrary geometries
\cite{Schafer2015}. That makes it natural to ask shape questions: where the
response concentrates, how it shifts under geometry perturbations, how much of
the signed weight lies in a chosen region, and whether the strongest response
tracks a designed boundary feature.

The White sphere--cylinder example is boundary-conditioned at this level. A
\(1\um\) diameter sphere centered in a nominal \(4\um\) diameter cylinder
produces a shell-like distribution in the reported worldline analysis
\cite{White2021}. The field-facing boundary question is whether a finite
sphere--wall geometry organizes the vacuum response into an annular shell
role. That question is independent from the larger source-scale and direct
readout questions. It asks what the boundary does.

\section{State-Conditioned Response}

The state-conditioned modality begins with a quantum field state rather than a
static boundary geometry. The signed object is a local or smeared
stress-energy feature of that state. Examples include squeezed-state
stress-energy profiles, measurement-conditioned field states, and
quantum-energy-teleportation-type protocols in which local operations and
classical records create a negative stress-energy density in a target region
\cite{Ford2009,FunaiMartinMartinez2017}.

The preparation record is part of the claim. It tells the reader how the field
state was made, which degrees of freedom were measured or driven, how the
classical record enters, where compensating positive energy appears, and what
time window carries the deficit. This modality is dynamic and operational. Its
strength is that the sign can be tied to a controllable state-preparation
procedure alongside fixed-boundary inference.

State-conditioned negative energy also naturally meets quantum energy
inequalities. Ford's review emphasizes the basic QFT fact that local negative
energy densities occur, with magnitude and duration governed by quantum
constraints \cite{Ford2009}. Kontou and Sanders place that fact inside the
broader energy condition landscape: quantum fields motivate a refined move
from pointwise conditions to averaged and quantum-inequality statements
\cite{KontouSanders2020}. For state-conditioned claims, duration, smearing,
compensation, and averaging are part of the observable content.

This modality is useful beside boundary-conditioned Casimir work because it
clarifies an important distinction. A Casimir geometry changes the vacuum
state through boundary conditions. A QET-style or squeezed-state protocol
changes, selects, or conditions the field state through operations. Both can
produce negative stress-energy language. Their preparation logic, compensation
logic, and readout logic are different.

\section{Tensor and Averaging Relevance}

The tensor/averaging modality asks which stress-energy object is being used.
For a relativistic source claim, the relevant object is more than a visual
negative region. It is a component or contraction of the renormalized stress
tensor, such as \(\Ttt\), \(\Tnn\), or an integral of \(\Tabmunu\) over a
specified sampling domain. The tensor object then has to be compared with the
source role under discussion.

This modality supplies the bridge from laboratory QFT to gravity-facing
language. Energy conditions in general relativity were historically formulated
as restrictions on stress-energy; quantum field theory gives well-understood
violations of pointwise versions and motivates averaged or quantum-inequality
replacements \cite{KontouSanders2020}. A laboratory negative-energy
measurement can therefore be precise about its level:
\begin{itemize}
\item local energy density or pressure in a material/boundary setup;
\item a stress-tensor component or finite-difference proxy;
\item a null or timelike contraction;
\item a smeared or averaged quantity;
\item a source-demand comparison for a specified metric model.
\end{itemize}

This layer is where morphology and source relevance separate cleanly. A
boundary-conditioned morphology can be real and reproducible while its
calibrated stress scale remains Casimir-scale relative to a gravity-facing
source demand. A state-conditioned deficit can be operationally generated
while its averaged integrals obey quantum-inequality restrictions. A readout can recover a
geometry-locked coefficient while the corresponding tensor component remains
an ordinary Casimir-scale object.

The positive role of the tensor/averaging modality is to keep each of those
statements useful. It lets morphology remain morphology, readout remain
readout, and source comparison remain source comparison. It gives the field a
way to promote claims by adding tensor and averaging evidence while preserving
the natural conclusion of each modality.

\section{Readout-Conditioned Observability}

The readout-conditioned modality begins with an instrument. Casimir
experiments measure force, force gradient, pressure, displacement,
oscillator-frequency shift, phase, quality factor, transition behavior, or
other apparatus responses. Those responses are sensitive to vacuum-boundary
physics through a transduction chain, and the same chain also carries ordinary
material and electromagnetic effects.

This modality is where laboratory claims become data statements. A readout
claim should identify:
\begin{itemize}
\item the observable and units;
\item the geometry schedule or state-preparation schedule;
\item calibration channels and null controls;
\item material, patch, thermal, vibration, alignment, and readout-chain
templates;
\item the recovery statistic and false-positive test.
\end{itemize}

Modern Casimir metrology already works in this style. Bimonte et al. combine
micromechanical torsional-oscillator measurements with explicit roughness,
edge, patch-potential, and optical-data modeling over \(0.2\)--\(8\um\)
separations \cite{Bimonte2021}. Patch-potential modeling and Kelvin-probe
measurements are central to high-quality force interpretation
\cite{Behunin2011,Garrett2015}. Superconducting-transition Casimir sensing
adds a cryogenic readout frontier in which tiny pressure changes are compared
against electrostatic, thermal, material, and alignment effects
\cite{Xu2026}.

Readout-conditioned observability is therefore a positive experimental
language. It gives a boundary or state result an instrument-native test. It
also keeps the measured object specific: pressure, frequency shift, phase,
loss, current, and voltage each carry a distinct transfer function, nuisance
structure, and claim boundary.

\section{Transfers Between Modalities}

Most strong laboratory papers move across modalities. A boundary calculation
motivates an instrument. A state-preparation protocol predicts a stress-tensor
deficit. A readout experiment recovers a coefficient that can be compared with
a tensor model. The transfer is scientifically productive when the promoted
claim carries the extra evidence required by the new modality.

\Tab~\ref{tab:transfers} lists common transfers and their evidence
requirements.

\begin{table}[h]
\centering
\small
\begin{tabularx}{\textwidth}{p{0.25\textwidth}p{0.31\textwidth}X}
\toprule
Transfer & Added evidence & Promoted conclusion \\
\midrule
Boundary morphology to tensor object &
Stress-tensor component, normalization, convergence, material model, and
smearing convention &
The shape is associated with a quantified energy, pressure, or stress channel. \\
\addlinespace
Tensor object to source comparison &
Metric source demand, units, sign convention, averaging window, and scale
ratio &
The laboratory stress-energy object can be compared with a declared
gravity-facing source role. \\
\addlinespace
Boundary or state object to readout &
Observable model, transfer function, controls, nuisance templates, and
false-positive test &
The apparatus can search for the corresponding coefficient in measured data. \\
\addlinespace
Readout coefficient to field interpretation &
Calibration closure, background separation, phase/geometry locking, and
replication under controls &
The measured signal can be assigned to the boundary or state-conditioned field
object under the declared model. \\
\addlinespace
State preparation to averaged-energy statement &
Sampling function, duration, compensating positive energy, and quantum
inequality comparison &
The local deficit is placed inside a finite-time and finite-region energy
ledger. \\
\bottomrule
\end{tabularx}
\caption{Evidence required when a claim moves from one modality to another.
The promoted conclusion names the new physical object rather than relabeling
the old one.}
\label{tab:transfers}
\end{table}

This transfer table is the practical core of the modality grammar. It
protects useful intermediate results. A morphology result can stand as a
morphology result. A readout architecture can stand as a readout architecture.
A source comparison can stand as a source comparison. The strongest papers
then show how those results connect.

\section{White Sphere--Cylinder Case Study}

The White sphere--cylinder result is an ideal case study because it touches
all four modalities. The original paper reports a boundary-conditioned
worldline-numerics result: a custom Casimir geometry produces a shell-like
negative distribution in a toy model consisting of a \(1\um\) sphere centered
in a \(4\um\) cylinder \cite{White2021}. The visual resemblance to an
Alcubierre-style source shape makes the tensor/source modality relevant. The
proposed voltage, current, photon, or electron detection ideas make the
readout modality relevant. The public interpretation then depends on keeping
those modalities aligned.

The companion audit in this paper ensemble treats the White geometry as a
boundary-conditioned and tensor/source case. Its role-resolved question is:
what remains after separating the sphere, cylinder wall, annular shell,
central readout path, far field, material/background accounting, and source
scale? Within that structure, the shell-like boundary morphology can be
supported while direct gravity-facing scale and central timing/readout routes
are bounded by calibration and recovery gates.

The companion readout paper treats the shell-coded object as a
readout-conditioned target. Its question is: which physical apparatus
ensembles could test a geometry-locked boundary response? Pressure cells,
high-\(Q\) cavities, and superconducting resonators become branch-native
readouts. Their success criteria are pressure contrast, eigenmode shift,
phase/loss response, or cryogenic microwave response with material and
electromagnetic controls.

Placed in the modality grammar, the case study has a clean identity:
\begin{itemize}
\item Boundary modality: the finite sphere--cylinder geometry organizes a
shell-like Casimir response.
\item Tensor/averaging modality: stress scale, source demand, and averaging
decide the gravity-facing interpretation.
\item Readout modality: pressure, cavity, superconducting, voltage, current,
photon, or electron channels decide what an apparatus can recover.
\item State modality: the White setup is primarily boundary-conditioned; any
state-conditioned extension would require an additional preparation protocol.
\end{itemize}

This is the value of the central paper. It lets the White pair stay
technically sharp while the ensemble has a shared conceptual grammar. It also
lets other negative-energy experiments locate themselves on the same map.

\section{Field-Facing Claim Grammar}

The field-facing version of the framework is a short checklist. A laboratory
negative-energy paper can state:
\begin{enumerate}
\item \textbf{Modality.} Is the sign boundary-conditioned, state-conditioned,
tensor/averaging-conditioned, readout-conditioned, or a declared combination?
\item \textbf{Object.} Is the object energy, pressure, force, stress,
frequency shift, phase, loss, current, voltage, detector response, or a model
coefficient?
\item \textbf{Support.} Where and when does the signed object live, and how is
the support defined?
\item \textbf{Compensation.} Where does the associated positive energy,
material response, reset cost, or background response appear?
\item \textbf{Controls.} Which material, electromagnetic, thermal, patch,
alignment, drift, and readout-chain controls share the same apparatus?
\item \textbf{Promotion path.} Which additional evidence would move the result
into a stronger modality?
\end{enumerate}

\Tab~\ref{tab:claimlevels} summarizes the resulting claim levels.

\begin{table}[h]
\centering
\small
\begin{tabularx}{\textwidth}{p{0.21\textwidth}p{0.31\textwidth}X}
\toprule
Claim level & Evidence basis & Clean conclusion \\
\midrule
Geometry organizes response &
Finite-boundary calculation, perturbation checks, morphology metrics &
The apparatus geometry carries a reproducible boundary-conditioned pattern. \\
\addlinespace
Stress channel quantified &
Stress tensor or calibrated proxy, material model, scale estimate &
The pattern has a quantified energy, pressure, or stress role. \\
\addlinespace
Operational state deficit &
Preparation record, detector/work ledger, duration and compensation &
The field state carries a controlled local or smeared deficit. \\
\addlinespace
Instrument recovers coefficient &
Measured observable, controls, template separation, false-positive test &
The apparatus sees a recoverable boundary or state-conditioned coefficient. \\
\addlinespace
Source relevance established &
Tensor component, averaging, source-demand comparison, scale closure &
The stress-energy object can be assigned to the declared gravity-facing source
role. \\
\bottomrule
\end{tabularx}
\caption{Claim levels enabled by the modality grammar. Each level is useful
on its own and can be promoted by adding the evidence named in the transfer
table.}
\label{tab:claimlevels}
\end{table}

This grammar also improves peer review. A reviewer can ask whether the paper's
evidence matches the modality being claimed. The authors can then answer at
the right level: boundary morphology, field-state preparation, stress-tensor
scale, readout recovery, or source relevance.

\section{Implications for Laboratory Programs}

The modality grammar suggests a constructive path for laboratory QFT programs.
First, identify the modality with the strongest existing evidence. For many
Casimir platforms, that is boundary-conditioned response plus
readout-conditioned force or pressure metrology. For QET-style protocols, it
is state-conditioned preparation plus operational energy extraction or detector
response. For energy-condition studies, it is tensor/averaging structure plus
duration and smearing constraints.

Second, design the transfer deliberately. A finite-boundary morphology needs a
stress or readout model before it becomes an apparatus claim. A readout signal
needs template separation before it becomes a field claim. A local
stress-energy deficit needs smearing and compensation before it becomes an
averaged-energy statement. A gravity-facing source comparison needs calibrated
stress scale, tensor component, sign convention, and source-demand ratio.

Third, treat ordinary apparatus physics as part of the field experiment.
Material response, patch potentials, finite conductivity, thermal behavior,
roughness, mode mixing, vibration, alignment, and readout-chain effects are
the channels through which laboratory negative-energy claims become precise.
In high-quality Casimir work, those
channels already provide the calibration structure that makes the measurement
scientific \cite{Bimonte2021,Behunin2011,Garrett2015,Xu2026}.

This has a positive consequence for speculative-adjacent systems. A visually
suggestive negative-energy morphology can be valuable at an early
source-assignment stage. It can be a boundary-conditioned target, a
stress-channel hypothesis, a readout template, or a discrimination schedule.
Each role supports real progress when it is named.

\section{Reproducibility and Source Basis}

The cited source PDFs for this central paper are mirrored in the local
\path{white_casimir_intersection/sources} directory. That directory includes
the White et al. EPJ C paper, QFT negative-energy and energy-condition
references, the worldline stress-tensor reference, engineered negative
stress-energy via quantum energy teleportation, and the Casimir metrology,
patch-potential, and superconducting-transition references used for the
readout discussion.

This manuscript is a conceptual synthesis. Its reproducible content is the
claim grammar: the modality definitions,
transfer-evidence table, and claim-level table. The two companion papers supply
the worked White sphere--cylinder audit and the physical readout architecture
application.

\section{Conclusion}

Negative-energy language is most powerful when its modality is explicit.
Boundary-conditioned response, state-conditioned response, tensor/averaging
relevance, and readout-conditioned observability are different laboratory
objects. They can reinforce one another, and their interfaces define strong
experimental programs.

The White sphere--cylinder case illustrates the payoff. The shell-like Casimir
boundary morphology, the stress/source scale, and the proposed readout routes
are connected but distinct. The central paper supplies the vocabulary for that
separation. The companion audit applies it to the White geometry. The
companion readout paper applies it to pressure, cavity, and superconducting
measurement ensembles.

The field-level result is a clean claim grammar for laboratory QFT. A paper
can now say which kind of negative energy it has, which object carries the
sign, which controls make the readout meaningful, and which additional
evidence would promote the result into a stronger modality. That grammar keeps
boundary effects, prepared field states, stress-tensor constraints, and
instrument responses scientifically connected while preserving the specific
meaning of each.

\begingroup
\raggedright
\begin{thebibliography}{20}

\bibitem{White2021}
H. White, J. Vera, A. Han, A. R. Bruccoleri, and J. MacArthur,
``Worldline numerics applied to custom Casimir geometry generates
unanticipated intersection with Alcubierre warp metric,''
\emph{European Physical Journal C} \textbf{81}, 677 (2021).
\url{https://doi.org/10.1140/epjc/s10052-021-09484-z}.

\bibitem{Schafer2015}
M. Schafer, I. Huet, and H. Gies,
``Worldline Numerics for Energy-Momentum Tensors in Casimir Geometries,''
arXiv:1509.03509 (2015).
\url{https://arxiv.org/abs/1509.03509}.

\bibitem{Ford2009}
L. H. Ford,
``Negative Energy Densities in Quantum Field Theory,''
arXiv:0911.3597 (2009).
\url{https://arxiv.org/abs/0911.3597}.

\bibitem{KontouSanders2020}
E.-A. Kontou and K. Sanders,
``Energy conditions in general relativity and quantum field theory,''
\emph{Classical and Quantum Gravity} \textbf{37}, 193001 (2020).
\url{https://doi.org/10.1088/1361-6382/ab8fcf}.

\bibitem{FunaiMartinMartinez2017}
N. Funai and E. Martin-Martinez,
``Engineering negative stress-energy densities with quantum energy
teleportation,''
\emph{Physical Review D} \textbf{96}, 025014 (2017).
\url{https://doi.org/10.1103/PhysRevD.96.025014}.

\bibitem{Decca2011}
R. S. Decca, V. Aksyuk, and D. Lopez,
``Casimir force in micro and nano electro mechanical systems,''
in \emph{Casimir Physics}, Lecture Notes in Physics \textbf{834} (2011).
\url{https://tsapps.nist.gov/publication/get_pdf.cfm?pub_id=906575}.

\bibitem{GarciaSanchez2012}
D. Garcia-Sanchez, K. Y. Fong, H. Bhaskaran, S. Lamoreaux, and H. X. Tang,
``Casimir Force and In Situ Surface Potential Measurements on
Nanomembranes,''
\emph{Physical Review Letters} \textbf{109}, 027202 (2012).
\url{https://doi.org/10.1103/PhysRevLett.109.027202}.

\bibitem{Stange2019}
A. Stange, M. Imboden, J. Javor, L. K. Barrett, and D. J. Bishop,
``Building a Casimir metrology platform with a commercial MEMS sensor,''
\emph{Microsystems \& Nanoengineering} \textbf{5}, 14 (2019).
\url{https://doi.org/10.1038/s41378-019-0054-5}.

\bibitem{Bimonte2021}
G. Bimonte, B. Spreng, P. A. Maia Neto, G.-L. Ingold,
G. L. Klimchitskaya, V. M. Mostepanenko, and R. S. Decca,
``Measurement of the Casimir Force between 0.2 and 8 um: Experimental
Procedures and Comparison with Theory,''
\emph{Universe} \textbf{7}, 93 (2021).
\url{https://doi.org/10.3390/universe7040093}.

\bibitem{Behunin2011}
R. O. Behunin, F. Intravaia, D. A. R. Dalvit, P. A. Maia Neto, and
S. Reynaud,
``Modeling electrostatic patch effects in Casimir force measurements,''
arXiv:1108.1761 (2011).
\url{https://arxiv.org/abs/1108.1761}.

\bibitem{Garrett2015}
J. L. Garrett, D. Somers, and J. N. Munday,
``The effect of patch potentials in Casimir force measurements determined by
heterodyne Kelvin probe force microscopy,''
\emph{Journal of Physics: Condensed Matter} \textbf{27}, 214012 (2015).
\url{https://doi.org/10.1088/0953-8984/27/21/214012}.

\bibitem{Xu2026}
M. Xu, R. J. G. Elbertse, A. Keskekler, G. Bimonte, J. Lee, S. Otte, and
R. A. Norte,
``High-Resolution Casimir Force Sensing Across a Superconducting Transition,''
arXiv:2504.10579, version 2 (2026).
\url{https://arxiv.org/abs/2504.10579}.

\end{thebibliography}
\endgroup

\end{document}
